\documentclass[10pt,a4paper]{amsart}
\usepackage{amsmath,amssymb,amsthm,graphicx}
\usepackage{verbatim}
\newenvironment{reference}{\comment}{\endcomment}
%\newenvironment{reference}{}{}
\newenvironment{slogan}{\comment}{\endcomment}
\newenvironment{history}{\comment}{\endcomment}

% For commutative diagrams we use Xy-pic
\usepackage[all]{xy}

% We use 2cell for 2-commutative diagrams.
\xyoption{2cell}
% Theorem environments.
%
\theoremstyle{plain}
\newtheorem{theorem}[subsection]{Theorem}
\newtheorem{proposition}[subsection]{Proposition}
\newtheorem{lemma}[subsection]{Lemma}

\theoremstyle{definition}
\newtheorem{definition}[subsection]{Definition}
\newtheorem{example}[subsection]{Example}
\newtheorem{exercise}[subsection]{Exercise}
\newtheorem{situation}[subsection]{Situation}

\theoremstyle{remark}
\newtheorem{remark}[subsection]{Remark}
\newtheorem{remarks}[subsection]{Remarks}

\numberwithin{equation}{subsection}

% Macros
%
\def\lim{\mathop{\mathrm{lim}}\nolimits}
\def\colim{\mathop{\mathrm{colim}}\nolimits}
\def\Spec{\mathop{\mathrm{Spec}}}
\def\Hom{\mathop{\mathrm{Hom}}\nolimits}
\def\Ext{\mathop{\mathrm{Ext}}\nolimits}
\def\SheafHom{\mathop{\mathcal{H}\!\mathit{om}}\nolimits}
\def\SheafExt{\mathop{\mathcal{E}\!\mathit{xt}}\nolimits}
\def\Sch{\mathit{Sch}}
\def\Mor{\mathop{\mathrm{Mor}}\nolimits}
\def\Ob{\mathop{\mathrm{Ob}}\nolimits}
\def\Sh{\mathop{\mathit{Sh}}\nolimits}
\def\NL{\mathop{N\!L}\nolimits}
\def\CH{\mathop{\mathrm{CH}}\nolimits}
\def\proetale{{pro\text{-}\acute{e}tale}}
\def\etale{{\acute{e}tale}}
\def\QCoh{\mathit{QCoh}}
\def\Ker{\mathop{\mathrm{Ker}}}
\def\Im{\mathop{\mathrm{Im}}}
\def\Coker{\mathop{\mathrm{Coker}}}
\def\Coim{\mathop{\mathrm{Coim}}}

% Boxtimes
%
\DeclareMathSymbol{\boxtimes}{\mathbin}{AMSa}{"02}

%
% Macros for moduli stacks/spaces
%
\def\QCohstack{\mathcal{QC}\!\mathit{oh}}
\def\Cohstack{\mathcal{C}\!\mathit{oh}}
\def\Spacesstack{\mathcal{S}\!\mathit{paces}}
\def\Quotfunctor{\mathrm{Quot}}
\def\Hilbfunctor{\mathrm{Hilb}}
\def\Curvesstack{\mathcal{C}\!\mathit{urves}}
\def\Polarizedstack{\mathcal{P}\!\mathit{olarized}}
\def\Complexesstack{\mathcal{C}\!\mathit{omplexes}}
% \Pic is the operator that assigns to X its picard group, usage \Pic(X)
% \Picardstack_{X/B} denotes the Picard stack of X over B
% \Picardfunctor_{X/B} denotes the Picard functor of X over B
\def\Pic{\mathop{\mathrm{Pic}}\nolimits}
\def\Picardstack{\mathcal{P}\!\mathit{ic}}
\def\Picardfunctor{\mathrm{Pic}}
\def\Deformationcategory{\mathcal{D}\!\mathit{ef}}
\usepackage{cleveref}
\crefname{originalreference}{original reference}{original references}
\Crefname{originalreference}{Original reference}{Original references}
\makeatletter
\newlabel{situation-shriek}{{0F42}{source}{}{sourcepointer-1}{}}
\newlabel{situation-shriek@cref}{{[originalreference][0][]0F42}{[1][0][]source}{}{}{}}
\makeatother
\crefname{originalreference}{original reference}{original references}
\Crefname{originalreference}{Original reference}{Original references}
\makeatletter
\newlabel{remark-composition-lower-shriek}{{0G54}{source}{}{sourcepointer-2}{}}
\newlabel{remark-composition-lower-shriek@cref}{{[originalreference][0][]0G54}{[1][0][]source}{}{}{}}
\makeatother
\crefname{originalreference}{original reference}{original references}
\Crefname{originalreference}{Original reference}{Original references}
\makeatletter
\newlabel{categories-remark-pro-category-copresheaves}{{0G2X}{source}{}{sourcepointer-3}{}}
\newlabel{categories-remark-pro-category-copresheaves@cref}{{[originalreference][0][]0G2X}{[1][0][]source}{}{}{}}
\makeatother
\crefname{originalproposition}{original proposition}{original propositions}
\Crefname{originalproposition}{Original proposition}{Original propositions}
\makeatletter
\newlabel{proposition-duality-compactly-supported}{{0G51}{source}{}{sourcepointer-4}{}}
\newlabel{proposition-duality-compactly-supported@cref}{{[originalproposition][0][]0G51}{[1][0][]source}{}{}{}}
\makeatother
\crefname{originallemma}{original lemma}{original lemmas}
\Crefname{originallemma}{Original lemma}{Original lemmas}
\makeatletter
\newlabel{lemma-upper-shriek-composition}{{0ATX}{source}{}{sourcepointer-5}{}}
\newlabel{lemma-upper-shriek-composition@cref}{{[originallemma][0][]0ATX}{[1][0][]source}{}{}{}}
\makeatother
\title{The lower shriek functors compose}
\author{The Stacks Project Authors}
\date{stacks-project:duality.tex:local-snapshot-20261012; GNU FDL 1.2 or later}
\begin{document}
\maketitle
\noindent\textit{Editorial note.} All delivered mathematical text is the author's own, copied from the pinned Stacks Project chapter file duality.tex; no statement or proof step was rewritten, reordered or removed.\par\medskip
\noindent\textit{Editorial note.} Preamble provenance: the pinned chapter file opens with an input of the shared preamble and carries no preamble text of its own, so the project's single Stacks-wide preamble staged at sources/stacks-duality/preamble.tex is excerpted verbatim under the recipe fixed by LANE\_BRIEF 22.1, as three spans, lines 8-13 (verbatim package and the reference/slogan/history environments), lines 14-18 (Xy-pic setup) and lines 154-217 (the native theorem declarations and mathematical macros). Nothing was added, removed or redefined; the wrapper class and the amsmath/amssymb/amsthm/graphicx base packages are supplied by the assembler.\par\medskip
\noindent\textit{Editorial note.} The counted statement and proof contain no cross reference, no citation and no included file, so the delivered body is self contained.\par\medskip
\noindent\textit{Editorial note.} Canonical location of the counted proposition: Stacks tag 0G55 (sources/stacks-duality/pinned-original-source.tex line 8882).\par\medskip
\noindent\textit{Allowed standard prerequisites (editorial).}\begin{itemize}
\item Duality for compactly supported pushforward. Same chapter, Proposition labelled proposition-duality-compactly-supported, cited at proof lines 8933-8945. Relates the lower shriek pushforwards to the upper shrieks.
\item Composition of the upper shriek functors. Same chapter, Lemma labelled lemma-upper-shriek-composition, cited at proof lines 8939-8941. Gives f\textasciicircum{}!g\textasciicircum{}! = (g o f)\textasciicircum{}!.
\item The pro-category description of morphisms of lower shriek objects. Categories, Remark labelled categories-remark-pro-category-copresheaves, cited at proof lines 8922-8924, and Same chapter, Remark labelled remark-composition-lower-shriek, invoked in the statement. Reduces the claim to testing Hom into bounded coherent objects.
\end{itemize}\medskip
\ifdefined\hypertarget\hypertarget{sourcepointer-1}{}\fi
\noindent\textit{Source reference.} Original reference 0F42: sources/stacks-duality/pinned-original-source.tex, source lines 3953--3953. Same-chapter pointer into the author environment situation at source line 3953; invoked at proof line None of the counted proof. It is retained only as a navigation pointer; every load-bearing step is inside the delivered span.\par\medskip
\ifdefined\hypertarget\hypertarget{sourcepointer-2}{}\fi
\noindent\textit{Source reference.} Original reference 0G54: sources/stacks-duality/pinned-original-source.tex, source lines 8851--8851. Same-chapter pointer into the author environment remark at source line 8851; invoked at proof line 8895 of the counted proof. It is retained only as a navigation pointer; every load-bearing step is inside the delivered span.\par\medskip
\ifdefined\hypertarget\hypertarget{sourcepointer-3}{}\fi
\noindent\textit{Source reference.} Original reference 0G2X: sources/stacks-chow/tags.txt, source lines 20072--20072. Cross-chapter standard prerequisite located in the project tag list staged in the local snapshot; invoked at proof line 8891 of the counted proof. It is retained only as a navigation pointer; every load-bearing step is inside the delivered span.\par\medskip
\ifdefined\hypertarget\hypertarget{sourcepointer-4}{}\fi
\noindent\textit{Source reference.} Original proposition 0G51: sources/stacks-duality/pinned-original-source.tex, source lines 8721--8721. Same-chapter pointer into the author environment proposition at source line 8721; invoked at proof line 8912 of the counted proof. It is retained only as a navigation pointer; every load-bearing step is inside the delivered span.\par\medskip
\ifdefined\hypertarget\hypertarget{sourcepointer-5}{}\fi
\noindent\textit{Source reference.} Original lemma 0ATX: sources/stacks-duality/pinned-original-source.tex, source lines 4161--4161. Same-chapter pointer into the author environment lemma at source line 4161; invoked at proof line 8917 of the counted proof. It is retained only as a navigation pointer; every load-bearing step is inside the delivered span.\par\medskip
\begin{lemma}
\label{lemma-composition-lower-shriek}
In Situation \ref{situation-shriek} let $f : X \to Y$ and $g : Y \to Z$
be composable morphisms of $\textit{FTS}_S$. With notation as in
Remark \ref{remark-composition-lower-shriek} we have
$Rg_! \circ Rf_! = R(g \circ f)_!$.
\end{lemma}

\begin{proof}
By the discussion in
Categories, Remark \ref{categories-remark-pro-category-copresheaves}
it suffices to show that we obtain the same answer if we compute
$\Hom$ into $L$ in $D^b_{\textit{Coh}}(\mathcal{O}_Z)$. To do this
we compute, using the notation in
Remark \ref{remark-composition-lower-shriek}, as follows
\begin{align*}
\Hom_Z(Rg_!Rf_!K, L)
& =
\colim_n \Hom_Z(Rg_!M_n, L) \\
& =
\colim_n \Hom_Y(M_n, g^!L) \\
& =
\Hom_Y(Rf_!K, g^!L) \\
& =
\Hom_X(K, f^!g^!L) \\
& =
\Hom_X(K, (g \circ f)^!L) \\
& =
\Hom_Z(R(g \circ f)_!K, L)
\end{align*}
The first equality is the definition of $Rg_!Rf_!K$. The second
equality is Proposition \ref{proposition-duality-compactly-supported}
for $g$.
The third equality is the fact that $Rf_!K$ is given by $(M_n)$.
The fourth equality is
Proposition \ref{proposition-duality-compactly-supported} for $f$.
The fifth equality is Lemma \ref{lemma-upper-shriek-composition}.
The sixth is 
Proposition \ref{proposition-duality-compactly-supported} for $g \circ f$.
\end{proof}

\end{document}
